\documentclass[a4paper]{article}
\usepackage[a4paper, landscape, margin=1cm]{geometry}
\usepackage[utf8]{inputenc}
\usepackage[T1]{fontenc}
\usepackage{adjustbox}
\usepackage{tikz}
\usetikzlibrary{shapes, arrows, fit, calc, positioning}
\usepackage{tikz-uml}
\tikzumlset{fill class=white, text=black, draw=black}
\pagestyle{empty}

\begin{document}

\begin{center}
{\Large Diagramma delle classi -- HomeStock}

\vspace{0.4cm}
\begin{adjustbox}{max width=\textwidth, max totalheight=0.9\textheight}
\begin{tikzpicture}

% ================= UTENTI E CASA =================
\umlclass[x=0, y=0]{Utente}{
  Nome : string \\ Email : string \\ Password : string \\
  NotificheAttive : bool \\ GiorniPreavviso : int}{
  registrati(email : string, password : string) : void \\
  accedi(email : string, password : string) : bool \\
  creaCasa(nome : string) : Casa \\
  impostaNotifiche(attive : bool, giorni : int) : void \\
  attivaPremium() : void}

\umlclass[x=0, y=-10]{UtentePremium}{FineAbbonamento : date}{
  statisticheSpreco(da : date, a : date) : float}
\umlinherit{UtentePremium}{Utente}

\umlclass[x=12, y=0]{Casa}{Nome : string}{
  aggiungiPosizione(nome : string) : Posizione \\
  invitaMembro(email : string) : void \\
  cercaProdotti(testo : string, ordine : string) : Prodotto[*] \\
  prodottiInScadenza(giorni : int) : Prodotto[*] \\
  ricetteSuggerite() : Ricetta[*]}

\umlassoc[name=app, mult1=1..*, mult2=*]{Utente}{Casa}
\umlassocclass[x=6, y=-4.5]{Appartenenza}{app-1}{
  Ruolo : string \\ DataIngresso : date}{}

% ================= LISTA DELLA SPESA =================
\umlclass[x=23, y=0]{ListaSpesa}{}{
  aggiungi(nome : string, quantita : float) : VoceSpesa \\
  aggiungiMancanti(r : Ricetta) : void}
\umlcompo[mult1=1, mult2=1]{Casa}{ListaSpesa}

\umlclass[x=33, y=0]{VoceSpesa}{
  Nome : string \\ Quantita : float \\ Acquistata : bool}{
  segnaAcquistata() : void \\
  trasferisciIn(p : Posizione) : Prodotto}
\umlcompo[mult1=1, mult2=*]{ListaSpesa}{VoceSpesa}

% ================= POSIZIONI E PRODOTTI =================
\umlclass[x=12, y=-8]{Posizione}{Nome : string}{
  aggiungiProdotto(nome : string, quantita : float, scadenza : date) : Prodotto \\
  elencoProdotti() : Prodotto[*]}
\umlcompo[mult1=1, mult2=1..*]{Casa}{Posizione}

\umlclass[x=12, y=-15]{Prodotto}{
  Nome : string \\ Quantita : float \\ Unita : string \\
  Scadenza : date \\ Marca : string}{
  aumenta(q : float) : void \\
  diminuisci(q : float) : void \\
  spostaIn(p : Posizione) : void \\
  giorniAllaScadenza() : int}
\umlcompo[mult1=1, mult2=*]{Posizione}{Prodotto}

\umlclass[x=0, y=-16]{Categoria}{
  Nome : string \\ GiorniConservazione : int}{
  scadenzaSuggerita() : date}
\umlassoc[mult1=*, mult2=1]{Prodotto}{Categoria}

\umlclass[x=12, y=-24]{Movimento}{
  Tipo : string \\ Quantita : float \\ Data : date}{}
\umlcompo[mult1=1, mult2=*]{Prodotto}{Movimento}

\umlclass[x=0, y=-24]{SchedaProdotto}{
  CodiceABarre : string \\ Nome : string \\ Marca : string \\ Immagine : string}{
  cercaPerCodice(codice : string) : SchedaProdotto}
\umlassoc[mult1=*, mult2=0..1]{Prodotto}{SchedaProdotto}

% ================= RICETTE =================
\umlclass[x=23, y=-15]{Ingrediente}{
  Nome : string \\ Sinonimi : string[*] \\ DiBase : bool}{
  riconosce(p : Prodotto) : bool}
\umlassoc[mult1=*, mult2=0..1]{Prodotto}{Ingrediente}

\umlclass[x=33, y=-15]{Ricetta}{
  Titolo : string \\ Minuti : int \\ Porzioni : int \\ Passaggi : string[*]}{
  disponibilita(c : Casa) : float \\
  ingredientiMancanti(c : Casa) : Ingrediente[*] \\
  usaProdottiInScadenza(c : Casa) : bool}
\umlassoc[name=riga, mult1=1..*, mult2=*]{Ingrediente}{Ricetta}
\umlassocclass[x=28, y=-21]{RigaRicetta}{riga-1}{Quantita : string}{}

\end{tikzpicture}
\end{adjustbox}
\end{center}

\end{document}
